\section{Results: the diagnostic attribution archive}
\label{sec:results}

The primary diagnostic archive in this revision is a single-model hidden-cause attribution run over 24 constructed cases, three in each of the eight registered families, produced with the glm-5.2 model. The metrics below are recomputed from the raw per-case records; the sidecar report is checked against those rows and is refused if it disagrees.

Accuracy is 21/24 $= 0.875$, with the nominal Wilson score interval given in
Table~\ref{tab:P5-3}. The interval treats the 24 fixed constructed cases as
Bernoulli units; because they are not a probability sample and come from one
model run, it is not a population confidence interval. This is a descriptive
finite-suite diagnostic score, not protected fresh-task improvement, a matched
baseline comparison, or a test of the primary hypothesis.

Three residual errors are retained and are not relabeled as successes (Table~\ref{tab:P5-residual-errors}). Case P5-HC-002, whose protected cause was a retrieval miss, was attributed to a representation gap at medium confidence; case P5-HC-012, an environment, dependency or tool failure, was attributed to an implementation bug at high confidence; and case P5-HC-018, a representation gap, was attributed to a method-basis gap at high confidence.

Each error now has a distinct successor research identity. P5-RD-01 crosses
retrieval access with representation regime; P5-RD-02 crosses candidate bytes
with independently reconstructed environments; and P5-RD-03 crosses
representation expansion with method-class expansion. Their four cells retain
main effects, interaction, and joint-null outcomes separately. P5-RD-01 and
P5-RD-03 have not been executed. P5-RD-02 now has three immutable public-development
identities; none changes the diagnostic error or contributes corrected
attribution accuracy.

\idt{P5-RD-02.DEFECTS4J.LANG1.V1} ended \textsf{CANNOT\_CHECK} with
\idt{HARMFUL_OR_ADVERSE_CELL_PRESENT}. Its registered JDK~17 arm violated
Defects4J's Java~11 execution contract, and its raw-upstream-tree identity model
incorrectly rejected the prepared local commits and trees that Defects4J creates
during checkout. V2 was therefore preregistered as a materially new successor,
not used to rewrite V1. It replayed the public known fix for Defects4J Lang-1 as
one public development cluster across two content-addressed Java~11
distributions. All four V2 cells were valid: the buggy revision failed the
registered developer test in both environments and the known fixed revision
passed in both. Thus the registered contrasts were
\texttt{implementation\_effect}=1.0,
\texttt{environment\_effect}=0.0 and \texttt{interaction}=0, giving the primary
terminal \idt{IMPLEMENTATION_MAIN_EFFECT} with the separately retained harm
modifier \idt{NO_ADVERSE_CELL}.
These are exact registered contrasts within one source/bug/fix cluster, not an
estimated population main effect; the two Java distributions and four cells are
within-cluster conditions rather than independent scientific units.

A direct retention audit then found a provenance limitation in V2. Its runner
used \texttt{Path.with\_suffix} for both checkout and test phase stems, collapsing
the two phases to the same stdout/stderr paths in each cell. The retained files
validate all eight declared test-stream hashes, but none of the eight declared
checkout-stream hashes; those checkout bytes were overwritten and are not
recoverable from the V2 archive. This limitation does not rewrite the registered
V2 cell outcomes or terminal.

Because V2 outcomes were already known, \idt{P5-RD-02.DEFECTS4J.LANG1.V3} was frozen
as an audit-only archival successor rather than an independent replication. Its
four cells replayed the same local pattern, while all 16 phase-qualified stream
files matched their declared hashes and every cell retained four distinct phase
paths. Its \idt{ARCHIVAL_REPLAY_COMPLETE} terminal means only that this
successor preserved its own replay logs. It neither recovers the missing V2
checkout bytes nor adds confirmatory causal, generalization or protected
authority beyond its own archival-integrity result.

The scientific unit is one public Defects4J Lang-1 source/bug/fix cluster under
a local evaluator (cluster-level $n=1$), not the four cells, two Java
distributions, or sixteen V3 phase streams. It is a known-fix replay, not a generated
repair or a protected campaign, and establishes no population generality,
independent custody, external independence, promotion authority, protected
transfer or protected freshness. H1--H4 remain \textsf{CANNOT\_CHECK}, and the
bridge remains at
\idt{P5_BRIDGE_PUBLIC_DEVELOPMENT_BOUND__PROTECTED_PROVIDER_AND_SIX_ARM_EXECUTION_CANNOT_CHECK}.
The three successor studies still serve to test whether registered interventions
split the mixed revision fibres identified in
Section~\ref{sec:minimal-revision-theory}; if a protected pair survives every
registered intervention, the correct output remains \textsf{UNRESOLVED}.

\subsection*{Terminal-adapter interface result}

The V2 source-interface audit supplies a formal result, not a comparator
outcome. Six prospectively frozen same-visible-symptom pairs show that each raw
native status fibre contains at least two minimal revision classes; accordingly
zero raw native symptoms license a singleton class. Its deterministic synthetic
suite passes 90/90 contract cases, with 70 \texttt{UNRESOLVED} terminals and 20
certificate/write-surface proposed-action types. The V3 successor does not
relax that gate: it passes 231/231 declared fictional cases, with 40 records in
20 supported constant fibres and 191 \texttt{UNRESOLVED}; raw-native licences
remain zero and ScienceClaw supports no singleton. No comparator, native output
example, public or protected outcome, or performance table was used. The
conditional support counts (seven, four, four, four, one, and zero across the
six fixed arms) describe interface coverage only; the 231 rows are neither
independent units nor correct-repair observations. The execution ledger retains
108 blocker instances, 18 per arm. Zero of six remains confirmatory-ready, and
H1--H4 remain \textsf{CANNOT\_CHECK}.

The later outcome-blind C1 V4 preflight binds six SWE-agent execution fields,
making C1 9/21 bound and 12/21 blocking. It does not parse a scientific result:
authored synthetic native-terminal checks all remain
\texttt{UNRESOLVED} at the revision-class layer, and the isolation wrapper
refuses execution with all 12 blockers. The other five arms remain unchanged;
at this C1-only snapshot the merged panel has 102 blocker instances, zero ready
arms and no performance or superiority observation.

The distinct outcome-blind C2 MOSS V4 preflight binds four fields beyond its
three prior source identities, making C2 7/21 bound and 14/21 blocking. Its
schema-v6 parser preserves native status/verdict but emits
\texttt{UNRESOLVED}; exact host and 1,196-entry OpenClaw locks are retained.
The complete authoritative tree contains zero \texttt{benchmark/} paths even
though the native trial runner requires \texttt{benchmark/claw-eval}, so a full
native smoke is \textsc{CannotCheck}. Updating C2 changes the current panel
count from 102 to 98 blocker instances, not to a performance observation: zero
of six remains ready and H1--H4 remain \textsf{CANNOT\_CHECK}.

The distinct C3 DGM V4 preflight binds a closed native patch parser, no
fallbacks and a 21,600-second outer watchdog, making C3 6/21 bound and 15/21
blocking. Its 23 requirements declarations contain zero exact pins and are not
mislabelled as a resolved dependency lock. The synthetic
\texttt{PATCH\_CAPTURED} fixture retains native completion but maps to
\texttt{UNRESOLVED}; a protected-key fixture is refused. Replacing C3's 18 V3
blockers with 15 changes the current panel count from 98 to 95. Zero of six
remains ready and H1--H4 remain \textsf{CANNOT\_CHECK}; no native DGM, model,
task, benchmark or protected outcome was run.

The distinct C4 ADIAS V4 preflight binds a hashed report parser, an empty
fallback/search policy and a 21,600-second foreground watchdog, making C4 6/21
bound and 15/21 blocking. The 1,578-file source manifest is byte-addressed, but
four VCS dependencies are unpinned, the mutable host-networked container path
does not isolate provider keys, and bundled task rights lack component
licence/NOTICE files under \texttt{data/}. Replacing C4's 18 V3 blockers with
15 changes the current panel count from 95 to 92. Zero of six remains ready and
H1--H4 remain \textsf{CANNOT\_CHECK}; no native ADIAS model, task, benchmark or
protected outcome was run.

The distinct C5 Double Ratchet metric-only V4 preflight binds six fields beyond
the three source identities, making C5 9/21 bound and 12/21 blocking. Its
results-only write surface, resource envelope and 46-entry uv lock are
outcome-blind interface bindings. The official runner nevertheless regenerates
hosted solver outputs and reports its development-locked metric every round;
the tree contains no frozen split/result payload and no P5-native task surface.
Replacing C5's 18 V3 blockers with 12 changes the panel count from 92 to 86.
Zero of six remains ready and H1--H4 remain \textsf{CANNOT\_CHECK}; no native
model, benchmark, P5 case, protected panel or performance outcome was run.

The distinct C6 ScienceClaw V4 preflight binds two fields beyond its three
source identities, making C6 5/21 bound and 16/21 blocking. The strict
native-shaped draft parser preserves every fibre as \texttt{UNRESOLVED}, and
an outer process-group wallclock is bound. The released CLI provides no native
revision selector, nominal dry-run would execute an investigation, and the
model/tool/reasoning/skill/dependency fallbacks cannot all be disabled.
Replacing C6's 18 V3 blockers with 16 changes the panel count from 86 to 84.
Zero of six remains ready and H1--H4 remain \textsf{CANNOT\_CHECK}; no
ScienceClaw model, tool, service, benchmark, protected outcome or prior-result
payload was executed or decoded.

The subsequent panel-wide V5 audit changes no arm result. Direct recomputation
gives 42 bound and 84 blocking field instances out of 126, with zero of six
arms ready. Fourteen exact blocker equivalence classes collapse into five
evidence-producing work packages of sizes 25, 18, 18, 15 and eight. This is a
79-item reduction in the coordination queue (84 to five), not a reduction in
the 84 scientific blocker instances: V5 creates zero new field bindings and no
effect estimate. It also exposes an unmatched resource contract despite all
six wallclock fields being individually bound. Its outcome-blind local
validator passes 74 checks and re-verifies six parser digests, but those checks
authorize neither execution nor performance. Fresh packet assurance is not
uniform: C4's native validator rewrites a temporary-path-dependent audit hash,
so its frozen audit receipt is determined not to reproduce. The cleaned C6
packet lacks the \texttt{.source-audit} checkout required by its native
validator, so fresh C6 source validation remains \textsc{CannotCheck}. Both
adverse assurance findings are retained rather than suppressed.

The P5 V6 common-case successor is the first material closure produced by that
root registry. It freezes one complete rights-cleared Lang-1 candidate core and
six acceptance receipts, creating exactly two new bindings per arm: common
candidate-visible case bytes and common task/benchmark-content rights. The
panel therefore moves from 42/126 to 54/126 bound and from 84 to 72 blocking.
The result does not infer six arm environments from one source archive: all six
\texttt{runtime.task\_environment} fields remain blocking, zero arm is ready,
and H1--H4 remain \textsf{CANNOT\_CHECK}. The packet-local validator reports
741/741 checks and verifies 22 checksum entries without executing a model or
opening an outcome.

The V7 native-environment fan-out closes one further field without executing
an arm. Exact offline setup bytes and a fixed effective agent configuration
bind C1's \texttt{runtime.task\_environment}; the other five receipts remain
blocking because their arm-specific artifacts do not exist. The aggregate
census is therefore 55/126 bound and 71/126 blocking, with five R2 instances
and zero of six arms ready. The packet validator reports one environment bound,
five blocking, 67 byte references verified and 25 packet files hashed.

The C3 V8/V9 successors close no further field. V8 builds a lawful 69-member
candidate seed and excludes 1,595 prior-outcome files, but the unchanged native
initializer requires the excluded \texttt{initial/overall\_performance} state.
V9's one static exact-byte/AST gate proves that an outcome-free adapter would
either fabricate performance, restore excluded outcomes, eliminate the
nondegenerate self-edit action or replace native parent-selection semantics.
The C3 environment remains blocking, so the aggregate stays 55/126 bound,
71/126 blocking and 0/6 ready. No DGM, model, benchmark, scorer or outcome is
executed.

Three later C2 packets test distinct lawful-successor obligations and must not
be read as a cumulative version chain. V11 binds
\texttt{runtime.task\_environment} for the separately authored
\idt{C2_LAWFUL_NATIVE_BYTE_SUCCESSOR__ORION_V11}, yielding 8/21 bound
and 13 blocking on the C2 V4 count basis. Its six registered public hexadecimal
classification classes pass the successor gate, but the task and evaluator are
authored post-outcome development instruments rather than released MOSS or a
protected panel. V12 separately binds
\texttt{inputs.candidate\_visible\_case\_bytes} and
\texttt{rights.task\_and\_benchmark\_content} for
\idt{C2_SOURCE_NATIVE_VISIBLE_CORE_SUCCESSOR__ORION_V12}, yielding
9/21 bound and 12 blocking. V12 does not inherit V11's runtime field.

V13 separately closes only
\texttt{rights.container\_and\_generated\_artifacts} for
\idt{C2_RIGHTS_CLEARED_SCRATCH_IMAGE_SUCCESSOR__ORION_V13}, yielding
8/21 bound and 13 blocking. Its exported Linux/arm64 scratch layer contains
exactly nine regular files and no symlinks. The normalized frozen-rootfs,
exported-layer, file-level-rights and SPDX~2.3 file sets are identical; full
Apache-2.0/NOTICE/CC0 bytes and explicit generated-artifact retention,
disclosure and publication authority accompany the image. The frozen runtime
receipt records exact stdout \texttt{ORION V13 container pass\string\n}, exit 0, no
network, read-only rootfs, all capabilities dropped, no-new-privileges and an
empty container diff. The image was archived, removed from the daemon and the
ephemeral remote build environment was destroyed.

None of V11--V13 inherits either other state, and their field closures cannot
be unioned. Released MOSS remains 7/21 bound and 14 blocking, and the separate
successors do not alter the authoritative 55/126 panel census. Consequently
0/6 arms remain ready and H1--H4, performance and superiority remain
\textsf{CANNOT\_CHECK}. No MOSS model, protected case, evaluator, scorer or
performance outcome was executed by these packets.

\subsection*{Frozen revision-level panel: no registered terminal}

The additive revision-level protocol verified its bound subject, protected
suite, split, evaluator and baseline identities and then executed ten frozen
policies on 96 constructed confirmatory cases, with 48 cases in each recorded
half. FULL\_T7 achieved revision-label accuracy $0.125$ --- \textbf{12/96}
cases --- false-broad rate
$0$, harmful-regression rate $0$ and fresh-transfer success $0.125$. Generic
causal diagnosis achieved accuracy and fresh-transfer success $1.0$ with zero
harm, while the random-diagnostic control reached $0.833$ on both measures.
Direct self-edit had accuracy $0$, false-broad rate $1.0$ and harmful-regression
rate $0.354$. The independent decision- and execution-layer checks agreed with
the primary scorer.

None of the seven frozen decision terminals fired: the execution receipt
records \texttt{terminal\_under\_frozen\_rules =
NO\_TERMINAL\_UNDER\_FROZEN\_RULES}, and that is the registered outcome rather
than an absence of one. The panel registered two
expected diagnostics for every hypothesis while the subject policy observed
one discriminator per bounded session; its responsibility gate therefore kept
every case unresolved until all modeled discriminators were observed. The
frozen panel and rules are retained unchanged. Within this panel H3's
cannot-check blocking and H4's history-blind re-admission check passed, whereas
H2 was not confirmed because the subject remained at the safety floor. These
are bounded secondary observations only: the receipt sets
\texttt{grants\_scientific\_authority=false}, H1 has no registered terminal,
and no general superiority or peer-review-readiness claim follows.

\subsection*{Pre-registered V4 revival of the revision-level panel}

The paragraph above closes with a contract gap rather than a software fault: the
panel registered two expected diagnostics per hypothesis while the bounded
subject observed one discriminator per session. The V3 lever addressing that gap
was frozen in the V3 receipt of 2026-08-24, \emph{before} any V4 panel or subject
existed, and was executed as a confirmatory successor on 2026-08-27. The V3
record is unchanged: \texttt{NO\_TERMINAL\_UNDER\_FROZEN\_RULES} and its $0.125$
accuracy stand exactly as reported, and the frozen V3 panel and rules are not
edited.

The successor panel carries 180 protected cases whose hypothesis expectation sets
are completable within the bounded protocol, together with ambiguous
\textsc{unresolved} cases and 20 preservation-conflict cases whose correct
terminal requires the revision-gate blocking branch the V3 panel left
unexercised. The subject \texttt{FULL\_T7\_V4} is the V3 chain with exactly one
mechanism completion --- candidate-visible \texttt{preserve:} obligations are
projected into the assessment step as obligation states plus forbidden writes, so
a diagnosis-licensed repair whose write coordinate is preserved blocks with
\texttt{FORBIDDEN\_WRITE} and the policy refuses rather than promoting. All ten
V3 baseline arms are delegated verbatim to the unmodified V3 implementation, so
the parent arm is the frozen V3 subject running unchanged.

Eleven frozen policies ran after all bindings were verified. \texttt{FULL\_T7\_V4}
reached revision-label accuracy $1.000$ with $0$ false-broad revisions, $0$
authority violations, correct-\textsc{unresolved} rate $1.000$ and preservation
refusal $1.000$. The parent \texttt{FULL\_T7} and \textsc{generic causal
diagnosis} both reached $0.889$ with $20$ authority violations and preservation
refusal $0.000$; \textsc{random diagnostic} reached $0.550$; \textsc{direct
self-edit} reached $0.000$ with $180$ false-broad revisions and $75$ harms. Under
the frozen surjective first-match decision rules, rule order 4 fired:
\texttt{REVISION\_LEVEL\_DISCRIMINATION\_SUPPORTED}. The subject was the only arm
satisfying the full frozen performance content, and every strong arm except the
subject took exactly $20$ authority violations, all concentrated in the
preservation stratum.

Three properties keep this from being a panel gift. The parent runs the same
panel unchanged and fails only there. The preservation wiring costs nothing
measurable: diagnostic cost is $170.0$ for subject, parent and generic alike,
because the gate blocks rather than spending an extra probe. And the subject is
feedback-dependent rather than cached --- accuracy $1.000$ under normal feedback
against $0.222$ under absent, permuted or contradictory feedback and $0.678$
under random feedback, where $0.222$ is the always-unresolved-compatible floor,
so it fails shut to refusal and never to promotion. Full re-execution reproduced
all $11 \times 180$ decision digests.

The boundary is narrow and is stated by the execution itself. Fresh-transfer
success is $0.889$ for the subject and $0.889$ for the strongest comparator: the
result is noninferiority on that measure, not improvement, so whether governed
method evolution improves fresh development outcomes remains undetermined. The
receipt sets \texttt{grants\_scientific\_authority=false} throughout; the terminal
is a frozen-rule selection over a benchmark-local protected panel; the
live-provider lane remains blocked; and the panel is constructed rather than
naturalistic. The execution was performed within the producing research group and
its independent decision- and execution-layer evaluators are second
implementations by the same researcher, so it is reproducibility evidence and is
\emph{not} external-investigator verification.

A separate scorer-independent audit recomputed the metrics, the performance
content and the terminal from the frozen committed receipt, and recovered the
same terminal. It also \textbf{found a defect in the executable decision rules},
which is disclosed here rather than repaired. The false-broad clause was declared
as \emph{false-broad rate no worse than every non-floor policy} but implemented as
\emph{no worse than the maximum non-floor policy rate}. The defect is real and is
retained: post-outcome in-place repair was not permitted, and the frozen rule
bytes stand as executed. It is non-controlling for this outcome only because both
the declared and the executable semantics independently yield the same rule-4
terminal with the same unique satisfying arm --- which is a property of this
panel's numbers, not a defence of the implementation. The audit's own state is
therefore \textsc{bounded\_verified\_with\_noncontrolling\_executable\_rule\_defect},
and it records \texttt{grants\_live\_provider\_claim}, \texttt{grants\_external\_real\_world\_validity}
and \texttt{grants\_human\_peer\_review\_authority} all false. That audit was a
separate-model session within the same research group and is likewise
reproducibility evidence, not external-investigator verification.

\subsection*{Post-outcome public-suite instrument diagnosis}

A separate NR-01 instrument study replays the unchanged 24-case public
hidden-cause suite. Its V1-control arm exactly reproduces 21/24 with standard
macro-F1 $0.872619$, including errors P5-HC-002, P5-HC-012 and P5-HC-018. A
staged-extraction-plus-deterministic-diagnosis instrument returns 24/24 with
standard macro-F1 $1.0$ on that same suite. This is an instrument-stage
positive, not evidence that a model acquired new attribution capability: the
requested model was GLM-5.2, the recorded served model was GLM-5.3, and the
control comparison pins the instrument rather than a model version.

The protocol, rows, report and revival receipt first enter repository history
together, so independent pre-outcome registration chronology is
\textsc{CannotCheck}. The suite was already public and frozen, the rows carry
no signed provider or independent-custody identity, and no blind, fresh or
source-disjoint transfer was run. The admissible disposition is therefore
\idt{POST_OUTCOME_PUBLIC_SUITE_INSTRUMENT_DIAGNOSIS_ONLY}. The historical
21/24 diagnostic remains immutable; H1--H4, protected freshness, comparative
performance and superiority remain \textsf{CANNOT\_CHECK}.

% Generated from archived glm-5.2 attribution JSONL. Do not edit by hand.
\begin{table}[t]
  \centering
  \caption{Hidden-cause attribution confusion matrix from the archived glm-5.2 diagnostic run over a fixed constructed suite. Accuracy is 21/24; the nominal Wilson score interval [0.690, 0.957] treats the 24 cases as Bernoulli units and is not a population confidence interval. Three residual errors are retained and are not relabeled as successes. This is not a protected fresh-transfer or H1 result.}
  \label{tab:P5-3}
  \small
  \begin{tabular}{lcccccccc}
    \toprule
    gold / attr. & Retr & Rout & Impl & Env & Eval & Repr & Meas & Meth \\
    \midrule
    Retr & 2 & 0 & 0 & 0 & 0 & 1 & 0 & 0 \\
    Rout & 0 & 3 & 0 & 0 & 0 & 0 & 0 & 0 \\
    Impl & 0 & 0 & 3 & 0 & 0 & 0 & 0 & 0 \\
    Env & 0 & 0 & 1 & 2 & 0 & 0 & 0 & 0 \\
    Eval & 0 & 0 & 0 & 0 & 3 & 0 & 0 & 0 \\
    Repr & 0 & 0 & 0 & 0 & 0 & 2 & 0 & 1 \\
    Meas & 0 & 0 & 0 & 0 & 0 & 0 & 3 & 0 \\
    Meth & 0 & 0 & 0 & 0 & 0 & 0 & 0 & 3 \\
    \bottomrule
  \end{tabular}
  \par\smallskip
  \textit{Residual errors:} P5-HC-002: gold RETRIEVAL\_MISS $\rightarrow$ REPRESENTATION\_GAP (MEDIUM) 
P5-HC-012: gold ENVIRONMENT\_DEPENDENCY\_TOOL\_FAILURE $\rightarrow$ IMPLEMENTATION\_BUG (HIGH) 
P5-HC-018: gold REPRESENTATION\_GAP $\rightarrow$ METHOD\_BASIS\_GAP (HIGH).
\end{table}

% Generated from archived glm-5.2 attribution JSONL. Do not edit by hand.
\begin{table}[t]
  \centering
  \caption{Retained residual attribution errors from the archived glm-5.2 diagnostic run. These three cases remain incorrect; they are not successes.}
  \label{tab:P5-residual-errors}
  \small
  \begin{tabular}{lp{0.30\linewidth}p{0.26\linewidth}l}
    \toprule
    case & gold cause & attributed cause & confidence \\
    \midrule
    P5-HC-002 & retrieval miss & representation gap & medium \\
    P5-HC-012 & environment dependency or tool failure & implementation bug & high \\
    P5-HC-018 & representation gap & method basis gap & high \\
    \bottomrule
  \end{tabular}
\end{table}


Replay-versus-fresh scatter, longitudinal specialist regression, the improvement/integrity frontier, failure recurrence, cost to protected improvement, the matched baseline and ablation comparison, and the campaign record of harmful and null interventions all have no admissible rows in this archive. Each of the seven tables below therefore names the measurement it would carry and leaves its result undetermined; no number in them is imputed from the diagnostic accuracy above.

% Generated from archived glm-5.2 attribution JSONL. Do not edit by hand.
% Status: CANNOT_CHECK — the archived attribution run has no motivating/replay/fresh
% score pairs, round identities, or matched-sample arms. Numbers are not imputed
% from the 21/24 diagnostic accuracy.
\begin{table}[t]
  \centering
  \caption{Replay-vs-fresh scatter: no admissible rows in the archived glm-5.2 attribution run. The archive contains a single-model hidden-cause diagnostic over 24 cases and no campaign-level replay/fresh score pairs. This table will be populated when a live protected-transfer campaign produces matched replay and fresh-task measurements.}
  \label{tab:P5-2}
  \small
  \begin{tabular}{lcc}
    \toprule
    Round & Replay score & Fresh score \\
    \midrule
    \multicolumn{3}{c}{\texttt{CANNOT\_CHECK} --- no admissible data in archive} \\
    \bottomrule
  \end{tabular}
  \par\smallskip
  \textit{Empirical authority:} DESCRIPTIVE\_ONLY (diagnostic attribution). Not a protected fresh-transfer or H1 result.
\end{table}
% Generated from archived glm-5.2 attribution JSONL. Do not edit by hand.
% Status: CANNOT_CHECK — the archived attribution run has no round identities,
% longitudinal specialist-designation annotations, or matched baseline/ablation
% arms. Numbers are not imputed from the 21/24 diagnostic accuracy.
\begin{table}[t]
  \centering
  \caption{Longitudinal specialist regression: no admissible rows in the archived glm-5.2 attribution run. The archive contains a single-model hidden-cause diagnostic over 24 cases and no longitudinal specialist-designation data. This table will be populated when a live campaign produces round-by-round specialist scores with regression annotations.}
  \label{tab:P5-4}
  \small
  \begin{tabular}{lccc}
    \toprule
    Round & Specialist & Prior score & Current score \\
    \midrule
    \multicolumn{4}{c}{\texttt{CANNOT\_CHECK} --- no admissible data in archive} \\
    \bottomrule
  \end{tabular}
  \par\smallskip
  \textit{Empirical authority:} DESCRIPTIVE\_ONLY (diagnostic attribution). Not a protected fresh-transfer or H1 result.
\end{table}
% Generated from archived glm-5.2 attribution JSONL. Do not edit by hand.
% Status: CANNOT_CHECK — the archived attribution run has no improvement/integrity
% frontier pairs, no matched-score arms, and no round identities. Numbers are not
% imputed from the 21/24 diagnostic accuracy.
\begin{table}[t]
  \centering
  \caption{Improvement-integrity frontier: no admissible rows in the archived glm-5.2 attribution run. The archive contains a single-model hidden-cause diagnostic over 24 cases and no frontier-pair data. This table will be populated when a live campaign produces matched improvement and integrity scores across rounds.}
  \label{tab:P5-5}
  \small
  \begin{tabular}{lcc}
    \toprule
    Round & Improvement score & Integrity score \\
    \midrule
    \multicolumn{3}{c}{\texttt{CANNOT\_CHECK} --- no admissible data in archive} \\
    \bottomrule
  \end{tabular}
  \par\smallskip
  \textit{Empirical authority:} DESCRIPTIVE\_ONLY (diagnostic attribution). Not a protected fresh-transfer or H1 result.
\end{table}
% Generated from archived glm-5.2 attribution JSONL. Do not edit by hand.
% Status: CANNOT_CHECK — the archived attribution run has no recurrence annotations,
% round identities, or matched-score arms. Numbers are not imputed from the 21/24
% diagnostic accuracy.
\begin{table}[t]
  \centering
  \caption{Failure recurrence: no admissible rows in the archived glm-5.2 attribution run. The archive contains a single-model hidden-cause diagnostic over 24 cases and no recurrence annotations. This table will be populated when a live campaign records which previously resolved failures recur across rounds.}
  \label{tab:P5-6}
  \small
  \begin{tabular}{lccc}
    \toprule
    Original failure & Round resolved & Round recurred & Outcome \\
    \midrule
    \multicolumn{4}{c}{\texttt{CANNOT\_CHECK} --- no admissible data in archive} \\
    \bottomrule
  \end{tabular}
  \par\smallskip
  \textit{Empirical authority:} DESCRIPTIVE\_ONLY (diagnostic attribution). Not a protected fresh-transfer or H1 result.
\end{table}
% Generated from archived glm-5.2 attribution JSONL. Do not edit by hand.
% Status: CANNOT_CHECK — the archived attribution run has no protected-improvement
% cost annotations, round identities, or matched-score arms. Numbers are not imputed
% from the 21/24 diagnostic accuracy.
\begin{table}[t]
  \centering
  \caption{Cost to validated improvement: no admissible rows in the archived glm-5.2 attribution run. The archive contains a single-model hidden-cause diagnostic over 24 cases and no cost-to-improvement annotations. This table will be populated when a live campaign records the resource cost (e.g., API calls, rounds, wall time) required to reach validated improvement.}
  \label{tab:P5-7}
  \small
  \begin{tabular}{lccc}
    \toprule
    Improvement & Validation & Cost metric & Value \\
    \midrule
    \multicolumn{4}{c}{\texttt{CANNOT\_CHECK} --- no admissible data in archive} \\
    \bottomrule
  \end{tabular}
  \par\smallskip
  \textit{Empirical authority:} DESCRIPTIVE\_ONLY (diagnostic attribution). Not a protected fresh-transfer or H1 result.
\end{table}
% Generated from archived glm-5.2 attribution JSONL. Do not edit by hand.
% Status: CANNOT_CHECK — the archived attribution run has no matched baseline/ablation
% arms, round identities, or campaign-level outcomes. Numbers are not imputed from
% the 21/24 diagnostic accuracy.
\begin{table}[t]
  \centering
  \caption{Baseline and ablation transfer/integrity results: no admissible rows in the archived glm-5.2 attribution run. The archive contains a single-model hidden-cause diagnostic over 24 cases and no baseline/ablation arm data. This table will be populated when a live campaign produces matched baseline-vs-ablation transfer and integrity measurements.}
  \label{tab:P5-T2}
  \small
  \begin{tabular}{lccc}
    \toprule
    Condition & Transfer score & Integrity score & $\Delta$ vs.\ baseline \\
    \midrule
    \multicolumn{4}{c}{\texttt{CANNOT\_CHECK} --- no admissible data in archive} \\
    \bottomrule
  \end{tabular}
  \par\smallskip
  \textit{Empirical authority:} DESCRIPTIVE\_ONLY (diagnostic attribution). Not a protected fresh-transfer or H1 result.
\end{table}
% Generated from archived glm-5.2 attribution JSONL. Do not edit by hand.
% Status: CANNOT_CHECK — the archived attribution run has no campaign intervention
% annotations, outcome labels, or round identities. Numbers are not imputed from
% the 21/24 diagnostic accuracy.
\begin{table}[t]
  \centering
  \caption{Campaign harmful/null interventions: no admissible rows in the archived glm-5.2 attribution run. The archive contains a single-model hidden-cause diagnostic over 24 cases and no intervention-level outcome annotations. This table will be populated when a live campaign records which interventions were harmful, null, or beneficial relative to the incumbent.}
  \label{tab:P5-T3}
  \small
  \begin{tabular}{lccc}
    \toprule
    Intervention & Outcome & Effect size & Recovery \\
    \midrule
    \multicolumn{4}{c}{\texttt{CANNOT\_CHECK} --- no admissible data in archive} \\
    \bottomrule
  \end{tabular}
  \par\smallskip
  \textit{Empirical authority:} DESCRIPTIVE\_ONLY (diagnostic attribution). Not a protected fresh-transfer or H1 result.
\end{table}
